% ============================================================
%  KWT: KDL Web Token — Technical White Paper
%  Overleaf-compatible LaTeX (pdflatex)
%  Packages: all available on Overleaf's TeX Live distribution
% ============================================================
\documentclass[11pt, letterpaper]{article}

% --- Core layout ---
\usepackage[
  top=1in, bottom=1in,
  left=1.15in, right=1.15in
]{geometry}
\usepackage{microtype}           % better line-break and spacing
\usepackage{parskip}             % paragraph spacing instead of indent
\setlength{\parskip}{6pt}

% --- Fonts ---
\usepackage[T1]{fontenc}
\usepackage{lmodern}

% --- Color ---
\usepackage[dvipsnames,table]{xcolor}
\definecolor{kwtblue}{HTML}{1A3A5C}
\definecolor{kwtmid}{HTML}{2C5F8A}
\definecolor{kwtlight}{HTML}{EAF1F8}
\definecolor{kwttablehead}{HTML}{1A3A5C}
\definecolor{codebg}{HTML}{F4F6F8}
\definecolor{codered}{HTML}{AA3333}
\definecolor{codegreen}{HTML}{226633}
\definecolor{codenavy}{HTML}{1A3A5C}

% --- Hyperlinks ---
\usepackage[
  colorlinks=true,
  linkcolor=kwtblue,
  urlcolor=kwtmid,
  citecolor=kwtblue,
  pdftitle={KWT: KDL Web Token},
  pdfauthor={KWT Working Group},
  pdfsubject={Authentication Protocol White Paper}
]{hyperref}

% --- Section heading styles ---
\usepackage{titlesec}
\titleformat{\section}
  {\color{kwtblue}\large\bfseries}
  {\thesection}{1em}{}
  [\color{kwtblue}\titlerule]
\titleformat{\subsection}
  {\color{kwtblue}\normalsize\bfseries}
  {\thesubsection}{1em}{}
\titleformat{\subsubsection}
  {\color{kwtmid}\normalsize\bfseries}
  {\thesubsubsection}{1em}{}
\titlespacing{\section}{0pt}{18pt}{6pt}
\titlespacing{\subsection}{0pt}{12pt}{4pt}
\titlespacing{\subsubsection}{0pt}{8pt}{2pt}

% --- Header/footer ---
\usepackage{fancyhdr}
\pagestyle{fancy}
\fancyhf{}
\fancyhead[L]{\small\color{gray}KWT: KDL Web Token}
\fancyhead[R]{\small\color{gray}Technical White Paper — Draft 1.0}
\fancyfoot[C]{\small\color{gray}\thepage}
\renewcommand{\headrulewidth}{0.4pt}
\renewcommand{\footrulewidth}{0pt}

% --- Tables ---
\usepackage{booktabs}
\usepackage{tabularx}
\usepackage{array}
\usepackage{multirow}
% Column type: centered, fixed width, with word wrap
\newcolumntype{C}[1]{>{\centering\arraybackslash}p{#1}}
\newcolumntype{L}[1]{>{\raggedright\arraybackslash}p{#1}}

% --- Code listings ---
\usepackage{listings}
\lstset{
  backgroundcolor=\color{codebg},
  basicstyle=\ttfamily\small,
  breaklines=true,
  breakatwhitespace=false,
  frame=single,
  rulecolor=\color{gray!40},
  framesep=4pt,
  xleftmargin=8pt,
  xrightmargin=8pt,
  captionpos=b,
  aboveskip=6pt,
  belowskip=6pt,
  columns=fullflexible,
  keepspaces=true,
  showstringspaces=false,
  tabsize=2,
  commentstyle=\color{gray},
  keywordstyle=\color{kwtblue}\bfseries,
  stringstyle=\color{codegreen},
}
% KDL "language" definition (approximate)
\lstdefinelanguage{KDL}{
  morekeywords={subject, issued, expires, audience, roles, scopes, token, version, jti},
  morestring=[b]",
  morecomment=[l]{//},
  sensitive=false,
}

% --- Math ---
\usepackage{amsmath}
\usepackage{amsfonts}

% --- Misc ---
\usepackage{enumitem}
\usepackage{graphicx}
\usepackage{float}
\usepackage{caption}
\captionsetup{font=small, labelfont=bf}
\usepackage{tcolorbox}
\tcbuselibrary{skins, breakable}

% Highlight box for key equations
\newtcolorbox{eqbox}{
  colback=kwtlight,
  colframe=kwtblue,
  fontupper=\ttfamily\normalsize,
  boxrule=0.6pt,
  left=6pt, right=6pt, top=4pt, bottom=4pt
}

% Warning / note box
\newtcolorbox{notebox}[1][Note]{
  colback=yellow!8!white,
  colframe=orange!60!black,
  title=\textbf{#1},
  fonttitle=\small,
  fontupper=\small,
  boxrule=0.6pt,
  breakable
}

% ============================================================
\begin{document}

% ============================================================
%  TITLE PAGE
% ============================================================
\begin{titlepage}
\centering
\vspace*{2cm}

{\color{kwtblue}\Huge\bfseries KWT}\par
\vspace{0.4cm}
{\color{kwtmid}\LARGE KDL Web Token}\par
\vspace{0.6cm}
{\large A Compact, Hardened Authentication Protocol}\par
\vspace{1.2cm}

\noindent\rule{\linewidth}{0.6pt}\par
\vspace{0.4cm}
{\large Technical White Paper \quad|\quad Draft 1.0}\par
\vspace{0.4cm}
\noindent\rule{\linewidth}{0.6pt}\par

\vfill
{\small\color{gray} KWT Working Group \\ \today}
\end{titlepage}

% ============================================================
\tableofcontents
\newpage

% ============================================================
\section{Abstract}
% ============================================================

JSON Web Tokens (JWT) are the dominant authentication primitive on the web.
They are flexible, widely supported, and well-understood.
They are also verbose, structurally under-constrained, and carry a long history
of catastrophic implementation vulnerabilities rooted directly in that flexibility.

This paper introduces \textbf{KWT} (KDL Web Token), a protocol that replaces
JWT's JSON payload encoding with KDL (KDL Document Language), couples it with
a versioned, algorithm-fixed cryptographic envelope, and defines a canonical
binary serialization that preserves KDL's information density advantage through
the encryption boundary.

The central thesis is that \emph{information density} — defined as meaningful
bits of claim data per byte of wire transmission — is a measurable, compoundable
engineering property. KWT achieves 2.5--4$\times$ the information density of
JWT for equivalent payloads, translating directly into reduced storage costs,
lower transmission latency, less CPU time for cryptographic operations, and
lower energy expenditure at scale.

We derive the mathematical basis for these claims, specify the complete protocol,
and provide a production Rust implementation (in-tree cryptographic primitives).

% ============================================================
\section{Introduction and Motivation}
% ============================================================

\subsection{The JWT Landscape}

JWT (RFC 7519) has become ubiquitous in web authentication, API authorization,
and session management. The format encodes a header, payload, and signature
as three base64url-encoded segments separated by periods:

\begin{lstlisting}
base64url(header).base64url(payload).base64url(signature)
\end{lstlisting}

JWT's popularity is justified: the format is human-readable before encoding,
self-contained, and stateless. However, the specification's emphasis on
flexibility has produced a well-documented catalogue of vulnerabilities:

\begin{itemize}[noitemsep]
  \item The \texttt{alg:none} attack — removing the signature by declaring
        no algorithm is in use
  \item Algorithm confusion — confusing HMAC and RSA families to forge
        signatures with public keys
  \item \texttt{kid} header injection — passing a key identifier to a
        database query unsanitized
  \item Weak secret keys — brute-forcing \texttt{HS256} with a short secret
        offline against intercepted tokens
\end{itemize}

These are not implementation bugs. They are the natural consequence of a format
that lets the token itself negotiate its own security parameters — a
fundamentally unsafe design choice.

\subsection{The Size Problem}

Beyond security, JWT carries a structural inefficiency that compounds at scale.
A representative JWT for a typical web session:

\begin{lstlisting}[language={}]
{"alg":"RS256","typ":"JWT"}
.{"sub":"user_882","iat":1744459200,"exp":1744545600,
  "roles":["admin","editor"],"aud":"api.example.com"}
.MEUCIQDd3...{86 bytes of base64 signature}
\end{lstlisting}

The encoded payload alone is approximately 160--180 bytes. At 100 million
authentications per day — a modest figure for a production service — this
represents roughly 17\,GB of payload data transmitted daily, excluding headers
and signature. Every byte stored in session databases, written to audit logs,
and processed by middleware multiplies the cost.

\subsection{The KDL Opportunity}

KDL (KDL Document Language, pronounced ``Cuddle'') is a human-readable document
language designed to eliminate the structural noise of JSON while preserving
its semantic expressiveness. The same payload in KDL:

\begin{lstlisting}[language=KDL]
subject "user_882"
issued 1744459200
expires 1744545600
audience "api.example.com"
roles "admin" "editor"
\end{lstlisting}

This is 91 characters versus JSON's approximately 115 characters — a 21\%
raw reduction before any encoding. More importantly, KDL's structure maps
cleanly to a compact canonical binary representation, allowing the density
advantage to survive encryption.

% ============================================================
\section{Information Density: The Calculus}
% ============================================================

\subsection{Defining Information Density}

We define information density $D$ for a token format as the ratio of semantic
payload bits $S$ to total transmitted bytes $W$:

\begin{equation}
  D = \frac{S}{W}
  \label{eq:density}
\end{equation}

where $S$ is the count of bits required to represent the minimum lossless
encoding of all claim values (excluding field names, which are structural
overhead), and $W$ is the total wire byte count of the transmitted token.
A higher $D$ means more claim data per transmitted byte.

\subsection{Structural Overhead Analysis}

Consider a canonical benchmark payload with five claims:

\begin{table}[H]
\centering
\caption{Semantic bit content of benchmark payload}
\small
\begin{tabularx}{\linewidth}{L{2.4cm} L{3.4cm} L{2.4cm} X}
\toprule
\rowcolor{kwttablehead}
\color{white}\textbf{Claim} &
\color{white}\textbf{Value} &
\color{white}\textbf{Semantic bits} &
\color{white}\textbf{Encoding note} \\
\midrule
\texttt{subject}  & user\_882         & $\sim$40 bits  & Integer user ID space \\
\texttt{issued}   & 1744459200        & 32 bits        & Unix timestamp \\
\texttt{expires}  & 1744545600        & 32 bits        & Unix timestamp \\
\texttt{audience} & api.example.com   & $\sim$112 bits & UTF-8 string \\
\texttt{roles}    & [admin, editor]   & $\sim$16 bits  & 2-bit enum $\times$ 2 \\
\bottomrule
\end{tabularx}
\end{table}

Total semantic content: approximately \textbf{232 bits (29 bytes)}.
Everything else is structural overhead: field names, delimiters, encoding
artefacts.

\subsection{Per-Format Overhead Breakdown}

\subsubsection{JSON Encoding}

\begin{lstlisting}[language={}]
{"sub":"user_882","iat":1744459200,"exp":1744545600,
 "aud":"api.example.com","roles":["admin","editor"]}
\end{lstlisting}

Raw byte count: 112 bytes. After base64url encoding: \textbf{150 bytes}.
Structural overhead is 61 bytes (54\%).

\begin{eqbox}
D(JWT-JSON) = 232 \,\text{bits} \,/\, 150 \,\text{bytes} = 1.55 \,\text{bits/byte}
\end{eqbox}

\subsubsection{KDL Text Encoding}

Raw KDL: 91 bytes. After base64url encoding: \textbf{122 bytes}.

\begin{eqbox}
D(KWT-text) = 232 \,\text{bits} \,/\, 122 \,\text{bytes} = 1.90 \,\text{bits/byte}
\end{eqbox}

\subsubsection{KWT Canonical Binary Encoding}

KWT's key innovation is a canonical binary serialization mapping KDL structure
to 1-byte opcodes, varint-encoded lengths, and enum values. The benchmark
payload encodes to \textbf{43 bytes} before encryption. After
XChaCha20-Poly1305 (which appends only a fixed 16-byte authentication tag)
and base64url encoding: approximately \textbf{79 bytes} for the ciphertext
segment. Combined with nonce and version prefix, total token: $\sim$\textbf{105--120 bytes}.

\begin{eqbox}
D(KWT-binary) = 232 \,\text{bits} \,/\, 79 \,\text{bytes} = 2.94 \,\text{bits/byte}
\end{eqbox}

\subsection{Comparative Summary}

\begin{table}[H]
\centering
\caption{Information density comparison across token formats}
\small
\begin{tabularx}{\linewidth}{L{3.6cm} C{2.2cm} C{2.2cm} C{2.6cm}}
\toprule
\rowcolor{kwttablehead}
\color{white}\textbf{Format} &
\color{white}\textbf{Pre-encrypt} &
\color{white}\textbf{Full token} &
\color{white}\textbf{Density (bits/byte)} \\
\midrule
JWT (JSON, RS256)         & 112 B & 380--500 B & 1.55 \\
KWT (KDL text)            & 91 B  & 220--380 B & 1.90 \\
KWT (canonical binary)    & 43 B  & 105--160 B & \textbf{2.94} \\
Theoretical minimum       & 29 B  & $\sim$65 B  & 3.57 \\
\bottomrule
\end{tabularx}
\end{table}

The canonical binary encoding achieves 89\% of the theoretical
information-density ceiling while remaining fully self-describing and verifiable.

\subsection{The Compounding Cost Model}

Let $T$ be daily token volume, $B(f)$ the bytes per token for format $f$,
$C_s$ the cost per GB-month of storage, $C_t$ the cost per GB of transmission,
and $C_e$ the CPU energy cost per byte decrypted.

\begin{equation}
  \text{Daily savings} = T \cdot \bigl(B(\text{JWT}) - B(\text{KWT})\bigr)
                         \cdot (C_s + C_t + C_e)
  \label{eq:savings}
\end{equation}

For a service at $T = 10^8$ tokens/day with $B(\text{JWT}) = 400$ bytes
and $B(\text{KWT}) = 120$ bytes:

\begin{itemize}[noitemsep]
  \item Payload reduction: 280 bytes per token
  \item Daily data reduction: \textbf{28 GB/day}
  \item Annual data reduction: \textbf{$\sim$10.2 TB/year}
\end{itemize}

At AWS S3 standard pricing ($\approx$\$0.023/GB-month), this yields
approximately \$2{,}800/year in storage savings alone before transmission
and compute are considered.

% ============================================================
\section{KWT Protocol Specification}
% ============================================================

\subsection{Design Principles}

KWT is designed around three non-negotiable constraints:

\begin{enumerate}[noitemsep]
  \item The token never negotiates its own cryptographic parameters. The version
        prefix determines the algorithm; the payload cannot override it.
  \item The payload is always encrypted. There is no signed-but-unencrypted mode.
  \item The parser is maximally strict. Unknown fields, out-of-range values,
        and malformed structure are hard errors, never silent ignores.
\end{enumerate}

\subsection{Token Structure}

\begin{lstlisting}
v{N}.<base64url(nonce)>.<base64url(ciphertext_with_tag)>
\end{lstlisting}

There is no header segment. The version prefix \texttt{v\{N\}} \emph{is} the
header — it is not attacker-controlled, is not base64-encoded, and determines
the complete cryptographic profile.

\subsection{Version Table}

\begin{table}[H]
\centering
\caption{KWT version registry (permanent; no runtime negotiation)}
\small
\begin{tabularx}{\linewidth}{C{1.0cm} L{3.2cm} L{2.4cm} X}
\toprule
\rowcolor{kwttablehead}
\color{white}\textbf{Ver.} &
\color{white}\textbf{Encryption} &
\color{white}\textbf{KDF} &
\color{white}\textbf{Use Case} \\
\midrule
v1 & XChaCha20-Poly1305 & HKDF-SHA256 & Standard symmetric tokens \\
v2 & AES-256-GCM        & HKDF-SHA256 & Hardware-accelerated (AES-NI); \textbf{browser-native} \\
v3 & XChaCha20-Poly1305 & X25519 + HKDF & Asymmetric: public encrypt, private decrypt \\
\bottomrule
\end{tabularx}
\end{table}

\begin{notebox}[Browser note]
\texttt{v2} (AES-256-GCM) is the only version whose primitives are natively
available via the W3C Web Crypto API in all modern browsers. See
Section~\ref{sec:browser} for details.
\end{notebox}

\subsection{Optional streaming and chunked AEAD (non-default)}

The compact three-segment wire form (\texttt{v\{N\}.nonce.ciphertext}) is the
\textbf{primary} KWT profile: a bearer object sized and used like \textbf{JWT}
(single AEAD package, bounded length, typical \texttt{Authorization} contexts).

Some deployments may additionally require \textbf{confidentiality over large or
streaming bodies}. For that, a future specification may register an
\textbf{opt-in} profile with its own framing (e.g.\ fixed-size chunks, each
authenticated before release, explicit maximum chunk count and rekeying rules).
Such a profile \textbf{must not} replace the default path: validators and
gateways that only implement the compact form continue to treat KWT as a
small, single-shot token. Until that profile is specified, implementations
\textbf{must} use only the compact wire shape for JWT-style credentials.

\subsection{Canonical Binary Payload Format}

\subsubsection{Opcode Registry}

All registered claim fields are assigned a 1-byte opcode that encodes both
the field identity and the value type.

\begin{table}[H]
\centering
\caption{KWT opcode registry}
\small
\begin{tabularx}{\linewidth}{C{2.2cm} L{2.6cm} L{2.8cm} X}
\toprule
\rowcolor{kwttablehead}
\color{white}\textbf{Range} &
\color{white}\textbf{Claim} &
\color{white}\textbf{Wire encoding} &
\color{white}\textbf{Notes} \\
\midrule
\texttt{0x10--0x1F} & Subject identifiers & varint len + UTF-8 & \\
\texttt{0x20--0x2F} & Timestamps & uint32 LE (4 B) & Unix epoch, seconds \\
\texttt{0x30--0x3F} & Audience / issuer & varint len + UTF-8 & \\
\texttt{0x40--0x4F} & Roles & uint8 count + uint8[] & Closed enum \\
\texttt{0x50--0x5F} & Scopes & uint8 count + uint8[] & Closed enum \\
\texttt{0x60--0x6F} & JTI / nonce & 16 raw bytes & UUID v7 only \\
\texttt{0x70--0x7F} & Custom claims & varint len + bytes & App-defined \\
\texttt{0x80}       & END marker & none & Required, last byte \\
\bottomrule
\end{tabularx}
\end{table}

\subsubsection{Varint Encoding}

Variable-length integers use Protocol Buffers-style encoding: the low 7 bits
carry data; the high bit signals continuation. Values 0--127 encode in 1 byte;
128--16\,383 in 2 bytes. This is optimal for token field lengths.

\subsubsection{Canonicalization Requirements}

Before encryption, the binary payload \emph{must} be serialized in canonical
form:

\begin{itemize}[noitemsep]
  \item Fields appear in \emph{ascending opcode order}
  \item No duplicate opcodes
  \item Strings are valid UTF-8, NFC-normalized
  \item Timestamps: $\mathtt{expires\_at} > \mathtt{issued\_at}$
  \item The \texttt{0x80} END marker is the final byte
\end{itemize}

\subsection{Nonce Strategy}

XChaCha20-Poly1305 uses a 192-bit (24-byte) nonce. KWT generates nonces using
a CSPRNG. The probability of nonce collision over $2^{32}$ tokens under the
same key is approximately $2^{-128}$ — well below any practical threshold.

Counter-based nonces are explicitly disallowed unless the implementation can
guarantee strict monotonicity across distributed systems and process restarts.

\subsection{Key Derivation}

Raw symmetric keys are not used directly. KWT uses HKDF-SHA256:

\begin{align}
  K_{\text{session}} &= \mathrm{HKDF}\bigl(
    \text{IKM}=K_{\text{master}},\;
    \text{salt}=\text{nonce}[0{:}32],\;
    \text{info}=\texttt{"kwt-v1-encryption"}\bigr)
  \label{eq:hkdf}
\end{align}

This ensures that even a nonce collision cannot recover the master key.

\subsection{Token Validation Algorithm}

Validation steps must be performed in the order below.
Any failure returns a generic 401 — callers must not expose which step failed.

\begin{enumerate}[noitemsep]
  \item Parse version prefix; reject unknown versions immediately.
  \item Decode base64url segments; reject malformed base64.
  \item (Optional) Replay filter on wire token or JTI hash \emph{before} decryption
        to shed obviously replayed traffic without full AEAD cost.
  \item Derive session key via HKDF.
  \item Decrypt and authenticate (AEAD); reject if tag fails.
  \item Parse canonical binary payload; reject malformed opcodes.
  \item Reject if $\mathtt{issued\_at}$ exceeds validator wall time plus a bounded future
        skew (reference implementation: $\approx 60$\,s); accept slightly early client clocks.
  \item Reject if $\mathtt{expires\_at}$, plus a small leeway (reference: $\approx 30$\,s),
        is before $\mathtt{now}$.
  \item Run application JTI uniqueness check (e.g.\ \texttt{SETNX}); reject replayed \texttt{jti}.
  \item Validate \texttt{audience} matches server's registered audience.
\end{enumerate}

% ============================================================
\section{Security Analysis}
% ============================================================

\subsection{Attack Surface Comparison}

\begin{table}[H]
\centering
\caption{Attack surface comparison: JWT vs KWT}
\small
\begin{tabularx}{\linewidth}{L{3.2cm} L{5.0cm} X}
\toprule
\rowcolor{kwttablehead}
\color{white}\textbf{Attack Vector} &
\color{white}\textbf{JWT Exposure} &
\color{white}\textbf{KWT Exposure} \\
\midrule
\texttt{alg:none} bypass    & Critical — native to spec        & Impossible — no algorithm field \\
Algorithm confusion         & High — attacker controls header  & Impossible — version is immutable \\
\texttt{kid} injection      & High — arbitrary string to DB    & Mitigated — key registry by UUID \\
Payload disclosure          & High — only signed, not encrypted & None — always AEAD-encrypted \\
Weak secret brute-force     & Possible with short HS256 keys   & Mitigated by HKDF derivation \\
Token replay                & Optional (\texttt{jti})          & Required (bloom filter) \\
Timing side-channels        & Library-dependent                & Constant-time (RustCrypto) \\
\bottomrule
\end{tabularx}
\end{table}

\subsection{The Encrypted Payload Security Model}

Because the entire payload is encrypted, an attacker who intercepts a KWT
token learns nothing beyond the token's version and length. A JWT containing
personally identifiable information (user IDs, roles, email) that appears in
a server log or CDN access log constitutes a data exposure event. A KWT in
the same context discloses only a random-looking byte string.

\subsection{The No-Header Guarantee}

The version string \texttt{v1} is matched against a hard-coded list — it is
not parsed as a data structure. An attacker cannot construct a malformed header
to trigger unexpected parser behaviour: there is no header to parse.

\subsection{Bloom Filter Replay Prevention}

Stateless tokens are inherently replayable within their validity window.
KWT mandates replay prevention via a counting bloom filter in Redis:

\begin{itemize}[noitemsep]
  \item JTI checked on every validation request before decryption
  \item TTL on filter entries matches token max TTL
  \item False-positive rate of 0.1\% is acceptable (results in a rejected
        valid token, not an accepted replayed token)
  \item UUID v7's time-ordering allows efficient range-based eviction
\end{itemize}

% ============================================================
\section{Comparison with Existing Standards}
% ============================================================

\subsection{KWT vs.\ JWT}

\begin{table}[H]
\centering
\caption{KWT vs JWT feature comparison}
\small
\begin{tabularx}{\linewidth}{L{3.2cm} L{4.0cm} X}
\toprule
\rowcolor{kwttablehead}
\color{white}\textbf{Property} &
\color{white}\textbf{JWT} &
\color{white}\textbf{KWT} \\
\midrule
Payload format      & JSON (verbose)             & Binary canonical (compact) \\
Encryption          & Optional (JWE)             & Mandatory \\
Algorithm agility   & Full (dangerous)           & None (safe) \\
Header              & Attacker-controlled JSON   & Immutable version prefix \\
Typical token size  & 300--500 bytes             & 100--160 bytes \\
Info density        & $\sim$1.55 bits/byte       & $\sim$2.94 bits/byte \\
Library support     & Ubiquitous                 & Reference only (Rust) \\
Audit history       & Extensive                  & None (new protocol) \\
\bottomrule
\end{tabularx}
\end{table}

\subsection{KWT vs.\ PASETO}

PASETO solves the algorithm-agility problem identically: versioned tokens with
fixed cryptographic profiles. KWT is best understood as PASETO with a
binary-encoded KDL payload rather than JSON.

\begin{table}[H]
\centering
\caption{KWT vs PASETO feature comparison}
\small
\begin{tabularx}{\linewidth}{L{3.2cm} L{4.0cm} X}
\toprule
\rowcolor{kwttablehead}
\color{white}\textbf{Property} &
\color{white}\textbf{PASETO} &
\color{white}\textbf{KWT} \\
\midrule
Security model      & Versioned, fixed algos     & Versioned, fixed algos \\
Payload format      & JSON                       & Binary canonical \\
Typical payload     & 150--250 bytes             & 40--80 bytes \\
Info density        & $\sim$1.75 bits/byte       & $\sim$2.94 bits/byte \\
Library support     & 10+ languages              & Rust \texttt{kwt} crate (WASM-capable) \\
Audit status        & Multiple audits            & No audits \\
\bottomrule
\end{tabularx}
\end{table}

% ============================================================
\section{Browser Compatibility}
\label{sec:browser}
% ============================================================

\subsection{Web Crypto API Coverage}

The W3C Web Crypto API (\texttt{window.crypto.subtle}) provides the cryptographic
primitives available natively in browser JavaScript. Its algorithm support
determines which KWT versions can run without a userland library.

\begin{table}[H]
\centering
\caption{KWT primitive availability in the Web Crypto API}
\small
\begin{tabularx}{\linewidth}{L{3.6cm} C{1.8cm} X}
\toprule
\rowcolor{kwttablehead}
\color{white}\textbf{Primitive} &
\color{white}\textbf{Available} &
\color{white}\textbf{Notes} \\
\midrule
AES-256-GCM (v2)        & \textcolor{codegreen}{\textbf{Yes}} & Full support Chrome, Firefox, Safari, Edge \\
HKDF-SHA256             & \textcolor{codegreen}{\textbf{Yes}} & All major browsers \\
\texttt{getRandomValues} & \textcolor{codegreen}{\textbf{Yes}} & All browsers; use for nonce generation \\
XChaCha20-Poly1305 (v1) & \textcolor{codered}{\textbf{No}}  & Not in spec; open W3C issue \#223 since 2019 \\
ChaCha20-Poly1305 (96-bit nonce) & \textcolor{codered}{\textbf{No}} & Used in TLS 1.3 internally but not exposed \\
\bottomrule
\end{tabularx}
\end{table}

The practical consequence: \textbf{KWT v2 (AES-256-GCM) can be implemented
entirely with Web Crypto API primitives}. KWT v1 and v3 require either a
WebAssembly module or a JavaScript fallback library.

\subsection{Recommended Browser Strategy}

\begin{enumerate}
  \item \textbf{Use v2 (AES-256-GCM) for all browser-issued or browser-validated
        tokens.} AES-256-GCM with hardware AES-NI acceleration is cryptographically
        equivalent to XChaCha20-Poly1305 for this use case; the nonce-reuse risk
        is managed by always generating nonces via \texttt{crypto.getRandomValues()}.

  \item \textbf{If v1 is required} (e.g., for interoperability with a Rust
        backend that issued v1 tokens), ship a \texttt{libsodium-wasm} bundle
        ($\sim$163\,kB gzipped) or compile the Rust \texttt{kwt} crate
        to WASM via \texttt{wasm-pack}.

  \item \textbf{Never place the master key in browser storage.}
        Browsers should receive a \emph{derived session key} or a short-lived
        \emph{wrapped key} via a secure server handshake — never the raw
        \texttt{MasterKey} bytes.
\end{enumerate}

\subsection{v2 JavaScript Reference (Web Crypto API)}

\begin{lstlisting}[language=JavaScript, caption={KWT v2 token validation in browser JavaScript}]
// KWT v2 browser validator (AES-256-GCM + HKDF-SHA256)
// The rawKey should arrive via a secure key-agreement handshake,
// never hardcoded or stored in localStorage.

async function kwtValidate(tokenStr, rawKey, expectedAudience) {
  const parts = tokenStr.split('.');
  if (parts.length !== 3 || parts[0] !== 'v2') {
    throw new Error('malformed token or wrong version');
  }

  const nonce = base64urlDecode(parts[1]);  // 12 bytes for AES-GCM
  const ciphertext = base64urlDecode(parts[2]);

  // 1. Import master key material
  const keyMaterial = await crypto.subtle.importKey(
    'raw', rawKey, 'HKDF', false, ['deriveKey']
  );

  // 2. Derive session key via HKDF-SHA256
  const sessionKey = await crypto.subtle.deriveKey(
    {
      name: 'HKDF',
      hash: 'SHA-256',
      salt: nonce.slice(0, 32).padEnd(32, 0), // pad 12-byte nonce to 32
      info: new TextEncoder().encode('kwt-v2-encryption'),
    },
    keyMaterial,
    { name: 'AES-GCM', length: 256 },
    false,
    ['decrypt']
  );

  // 3. Decrypt and authenticate
  const plaintext = await crypto.subtle.decrypt(
    { name: 'AES-GCM', iv: nonce },
    sessionKey,
    ciphertext
  );

  // 4. Parse canonical binary payload
  const claims = parseKwtBinary(new Uint8Array(plaintext));

  // 5. Validate expiry (with 30s leeway)
  const now = Math.floor(Date.now() / 1000);
  if (claims.expiresAt + 30 < now) throw new Error('token expired');

  // 6. Validate audience
  if (claims.audience !== expectedAudience) throw new Error('audience mismatch');

  return claims;
}
\end{lstlisting}

% ============================================================
\section{Implementation Guidance}
% ============================================================

\subsection{Rust implementation components}

\begin{table}[H]
\centering
\caption{Cryptography and dependencies in the production \texttt{kwt} crate}
\small
\begin{tabularx}{\linewidth}{L{3.6cm} L{3.4cm} X}
\toprule
\rowcolor{kwttablehead}
\color{white}\textbf{Component} &
\color{white}\textbf{Purpose} &
\color{white}\textbf{Notes} \\
\midrule
\texttt{primitive::xchacha20poly1305} & XChaCha20-Poly1305 (v1) & In-tree; RFC 8439 + draft XChaCha \\
\texttt{primitive::hkdf} + \texttt{primitive::sha256} & HKDF-SHA256 & In-tree; RFC 5869 \\
\texttt{getrandom}        & OS CSPRNG                 & Nonces and master-key generation \\
\texttt{uuid}             & UUID v7 for JTI           & Time-ordered for replay indexing \\
\texttt{zeroize}          & Key material erasure      & Clears sensitive buffers on drop \\
\texttt{base64ct}         & base64url wire encoding   & Constant-time decode \\
\bottomrule
\end{tabularx}
\end{table}

\subsection{Critical Implementation Notes}

\begin{itemize}
  \item \textbf{Zeroize key material.} Rust does not guarantee memory erasure
        on drop. Use the \texttt{zeroize} crate explicitly on all key buffers.
  \item \textbf{Fuzz-test the canonical codec.} Serialize-deserialize round-trip
        fuzzing with \texttt{cargo-fuzz} must be completed before deployment.
  \item \textbf{Bloom filter eviction must be load-tested.} An unbounded bloom
        filter will eventually produce an unacceptable false-positive rate.
  \item \textbf{Allow clock leeway.} Service-to-service clock skew can cause
        valid tokens to appear expired. A configurable leeway of 30--60 seconds
        is recommended.
  \item \textbf{Maintain validation order.} The bloom filter check must precede
        decryption to prevent replay-flood denial-of-service.
\end{itemize}

\subsection{Migration Path from JWT}

Migration proceeds in three phases:

\begin{enumerate}
  \item \textbf{Dual-accept.} The server accepts both JWT and KWT. New tokens
        are issued as KWT. Old JWT sessions expire naturally over the TTL window.
  \item \textbf{KWT-only.} After the dual-accept window, JWT acceptance is
        removed. All clients must be issuing and receiving KWT.
  \item \textbf{Cleanup.} JWT parsing code is removed; JWT secret keys are
        rotated out and destroyed.
\end{enumerate}

% ============================================================
\section{Conclusion}
% ============================================================

KWT demonstrates that token format is a measurable engineering variable with
direct cost implications, not merely a cosmetic choice.
By replacing JSON's verbose, delimiter-heavy encoding with KDL's node-structured
model and a compact canonical binary serialization, KWT achieves approximately
\textbf{2.94 bits of semantic information per transmitted byte} —
compared to JWT's 1.55 bits/byte — a 90\% improvement in information density.

This density advantage compounds across storage, transmission, and compute
dimensions. At scale, the savings are material. At the protocol level, the
versioned, algorithm-fixed design eliminates the entire class of
flexibility-based JWT attacks that have dominated web security incident
reports for a decade.

The primary limitation of KWT at this stage is the absence of broad
multi-language implementations and independent security audits of every
deployment. The recommendation is to treat this specification as a design target:
harden and fuzz the Rust \texttt{kwt} crate (or your port), obtain professional
audit coverage appropriate to your threat model, and deploy behind a feature
flag in a high-throughput non-critical path before promoting to primary
authentication infrastructure.

% ============================================================
\section*{References}
\addcontentsline{toc}{section}{References}
% ============================================================

\begin{enumerate}[label={[\arabic*]}]
  \item M. Jones, J. Bradley, N. Sakimura.
        \textit{RFC 7519: JSON Web Token (JWT)}. IETF, May 2015.
  \item S. Arciszewski.
        \textit{PASETO: Platform-Agnostic Security Tokens}, v4 spec, 2022.
  \item D. J. Bernstein.
        \textit{ChaCha, a variant of Salsa20}. SASC 2008.
  \item KDL Document Language Specification. \url{https://kdl.dev}, 2024.
  \item H. Krawczyk, P. Eronen.
        \textit{RFC 5869: HMAC-based Key Derivation Function (HKDF)}. IETF, 2010.
  \item W3C.
        \textit{Web Cryptography API}, W3C Recommendation.
        \url{https://www.w3.org/TR/WebCryptoAPI/}
  \item W3C WebCrypto GitHub Issue \#223.
        \textit{Please add support for chacha20-poly1305}.
        \url{https://github.com/w3c/webcrypto/issues/223}
  \item MDN Web Docs.
        \textit{SubtleCrypto.deriveKey()}.
        \url{https://developer.mozilla.org/en-US/docs/Web/API/SubtleCrypto/deriveKey}
\end{enumerate}

\end{document}
